\section{Workstream 10: the report}
\label{sec:report}

The report is worth 12 points. It follows the brief's section list in the brief's order, so the
completeness point can't be lost. The word budget is weighted by marks: data collection and
processing carry 5 points, so they get the most words.

\subsection{Cover and front matter (not counted)}

\begin{itemize}
  \item \textbf{Title}: ``Bitcoin Address Activity Explorer: Detecting Unusual Transactions with
    Machine Learning''. ``Unusual'', never ``fraud''.
  \item \textbf{Group identities}: group name HCMC\_COS30049.1\_4; group ID (``Group 4'' on the
    contribution form); the four names and student IDs; \textbf{the tutor's name}; the unit;
    the submission date.
  \item The word count and how it was counted (Section~\ref{sec:wordcount}).
  \item Table of contents, list of figures, list of tables.
\end{itemize}

\subsection{Section-by-section plan}

\begin{longtable}{@{}L{2.6cm}C{1.25cm}L{7.15cm}L{3.2cm}@{}}
\toprule
\textbf{Section} & \textbf{Words} & \textbf{Content} & \textbf{Figures and tables} \\
\midrule
\endhead
1 Introduction & 170 &
Motivation: anyone who handles cryptocurrency must decide which of many transactions deserve a
closer look, and an address's history is a long, unexplained list. Intended user: a student
researcher examining one supported address (as in A1); analysts doing triage benefit too. What A2
delivers: data, models and evaluation that rank each transaction by how unusual it is for
\emph{its own address}, with a reason, plus labelled-data experiments showing what such scores
can and cannot say. The limitation line, stated once. &
-- \\
2 Problem framing & 210 &
The challenge: deciding which transactions deserve attention without trustworthy labels.
Limits of existing solutions: block explorers show raw lists; fixed global rules ignore each
address's scale; supervised models trained on Elliptic and BitcoinHeist rely on sparse, noisy
labels and on anonymised or aggregated features that can't be computed for an arbitrary address
(Weber et al.\ 2019; Elmougy and Liu 2023). Why ML fits: interactions between features,
expected behaviour learned per address, structure found without labels, sequence context, and
supervised learning where labels exist. Four research questions: RQ1, the best stable detector
of planted patterns; RQ2, supervised performance on labelled data over time and on unseen
families; RQ3, how far unusualness and labelled illicit activity overlap; RQ4, whether a
from-scratch sequence model learns more than habit statistics. &
-- \\
3 Data collection & 410 &
The sources table, then three short subsections. Our own collection (DS1, DS2): API, paging,
pacing, the selection rule fixed before fetching, the verification sums, and the manifest.
Public labelled data (DS3, DS4): citation, access, licence, what the labels mean. Text corpora
(DS5--DS7). Challenges and responses, then the selection-bias statement (DR-02 condition 3). &
T1 datasets; T2 challenges \\
4 Data processing & 400 &
Cleaning steps with row counts; amount and direction; F1--F4 and the windows; the robust z;
scorability; the time split; the leakage controls in 3--4 sentences; the Elliptic join and
temporal split; BitcoinHeist binarisation, transforms, excluded leakage columns and sampling;
TxGPT tokenisation; the injection design. &
F1 pipeline; T3 cleaning log; T4 features \\
5 Data analysis & 370 &
The five findings of Section~\ref{sec:analysis}, each written as observation, then statistic,
then decision. &
F2 distributions; F3 in vs out; F4 PCA density; F5 Elliptic balance over time \\
6 Model selection & 400 &
The criteria C1--C6; the condensed candidates table; one paragraph each for clustering,
regression, classification, the A1 extras and the two language models, saying why each was
kept. The selection rule for the Explorer's method. &
T5 candidates \\
7.1 Technical implementation & 170 &
Libraries and why (pandas, scikit-learn, matplotlib, PyTorch for the language models); package
structure and CLI; configuration in one file; seeds and pinned environments; hardware and
measured run times; the prediction interface ready for A3. &
(architecture is F1) \\
7.2 Implementation evaluation & 560 &
The objectives table first. Then the DS1 comparison (recall by pattern, stability, sanity,
agreement) and the selected method; the labelled results and the label check; the language-model
results; unexpected outcomes; limitations; ethics in two sentences. &
T6 objectives; T7 DS1 comparison; T8 labelled results; T9 LM summary; F6 recall; F7 agreement; F8 PR curves and F1 per step; F9 TxGPT \\
8 Conclusion & 160 &
One answer per RQ; the method A3 will use and why; three limitations; future work (graph neural
networks on Elliptic, OPTICS, more addresses, the A3 user study). &
F10 example flagged transaction \\
\midrule
\textbf{Total} & \textbf{2,850} & 150 words of margin under the 3,000 limit. & \\
\bottomrule
\end{longtable}

\subsection{Appendices (not counted)}

\begin{enumerate}[label=\Alph*.]
  \item Data dictionary summary (the full one is in \code{DATA_CARD.md}).
  \item Hyperparameter searches: every grid and the chosen values.
  \item Full results: every model, split and seed.
  \item Language models: configs, training curves, the TxGPT scaling plot, sample outputs.
  \item Promises from A1 and how A2 met them (Section~\ref{sec:a1promises}), plus a summary of
    CR-06.
  \item Usability walkthrough results (A1 team action T7), only if the sessions happened.
  \item Reproducibility: hardware, library versions, measured run times.
\end{enumerate}

\subsection{Writing rules}

\begin{itemize}
  \item Short sentences in plain English. The present tense describes the system; the past tense
    describes what was done. Each acronym is defined once.
  \item One term per concept, used everywhere: \emph{unusual}, \emph{flag}, \emph{planted
    transaction}, \emph{fit period}, \emph{evaluation period}, \emph{address history}.
  \item No banned word describes the tool or DS1. ``Illicit'' and ``ransomware'' appear only as
    the labelled datasets' own classes.
  \item Every number comes from \fp{reports/tables/}, and a script checks the report's numbers
    against the CSVs before submission.
  \item Every figure and table is referenced in the text. Every caption says what to notice, not
    only what is shown.
  \item Claims about other work are cited. Claims about our own work point to a table or figure.
\end{itemize}

\subsection{Referencing}

Numbered Harvard, as in A1: an alphabetical, numbered list with author--date citations in the
text. The draft list is Appendix~\ref{app:refs}. Every entry is checked against its source before
use: authors, year, title, venue, DOI or URL, and the access date for web sources. A small script
lists citations without entries and entries never cited. Both lists must be empty.

\subsection{Word count}
\label{sec:wordcount}

\code{texcount} isn't installed, so \code{tools/wordcount.py} strips LaTeX from the body sections
(Introduction to Conclusion) and counts the words. The brief excludes the cover, contents,
bibliography and appendix, but says nothing about tables and captions.

\begin{team}
Ask the tutor, in the CR-06 message, whether tables and captions count towards the 3,000 words.
Until they answer, the plan keeps body text at or below 2,850 words and tables lean. The cover
states the counting rule used, for example ``2,840 words: body text excluding tables, captions,
references and appendices''.
\end{team}

\subsection{How the report gets produced}

\begin{enumerate}
  \item LaTeX, reusing the A1 preamble. A1's preamble loads \code{xltabular}, \code{pdflscape},
    \code{csquotes} and \code{seqsplit}, which the local TinyTeX doesn't have, so they are
    replaced with \code{longtable}, plain tables and \code{\textbackslash path} to make the build
    local and repeatable.
  \item Figures are \code{\textbackslash includegraphics} of the PDFs in
    \fp{reports/figures/}. Tables are \code{\textbackslash input} from LaTeX files generated next
    to the CSVs, so re-running the evaluation updates the report.
  \item \textbf{Writing order.} Data sections on 1 November, once the data facts are frozen;
    models and evaluation on 2 November; introduction, problem framing and conclusion last, on
    3 November, so they match the actual results.
  \item \textbf{Review on 5 November.} Cam Dan Bui (report integrator in A1) reads it against
    the rubric and the completeness checklist and proofreads. Comments are resolved the same day.
\end{enumerate}

\subsection{Completeness checklist}

\begin{center}
\small
\begin{tabularx}{\textwidth}{@{}Y C{1.4cm}@{}}
\toprule
\textbf{Required by the brief} & \textbf{Check} \\
\midrule
Project title & $\square$ \\
Group name and ID; student names and IDs; tutor's name & $\square$ \\
Introduction: motivation and intended user & $\square$ \\
Problem framing: challenge, limits of existing solutions, why ML & $\square$ \\
Data collection: sources, methods, challenges and how they were addressed & $\square$ \\
Data processing: missing values, normalisation, feature engineering, joining datasets & $\square$ \\
Model selection: criteria, algorithms considered, rationale & $\square$ \\
Technical implementation: libraries, high-level design, reasons for final decisions & $\square$ \\
Implementation evaluation: results against objectives, performance, unexpected outcomes & $\square$ \\
Conclusion: findings, insights, future work & $\square$ \\
Bibliography in Harvard style & $\square$ \\
Appendix & $\square$ \\
At most 3,000 words; file named after the group; PDF & $\square$ \\
\bottomrule
\end{tabularx}
\end{center}
